\documentclass[11pt]{article}
\usepackage[letterpaper,margin=1in]{geometry}
\usepackage{amsmath,amssymb,amsthm}

\newtheorem{theorem}{Theorem}
\newtheorem{lemma}{Lemma}
\newtheorem{corollary}{Corollary}
\newtheorem{definition}{Definition}
\newtheorem{proof}{Proof}

\title{Supplementary Material: Formal IIAR Framework}
\author{}

\begin{document}
\maketitle

\section{Formal Definitions}

\begin{definition}[Prompt Sensitivity Jacobian]
For a language model $\mathcal{M}_\theta$ with parameters $\theta$ and prompt $p$, let $h_\theta(p) \in \mathbb{R}^d$ be the final hidden state. The \emph{prompt sensitivity Jacobian} is
\[
J_\theta(p) = \nabla_p h_\theta(p) \in \mathbb{R}^{d \times |p|}
\]
measuring how infinitesimal changes to each prompt token affect the model's representations.
\end{definition}

\begin{definition}[Attention Redistribution Factor]
Let $\theta_B$ and $\theta_I = \theta_B + \Delta\theta$ be base and instruct parameters. The attention redistribution factor is
\[
\kappa(\theta_B, \theta_I, p) = \frac{\|J_{\theta_I}(p)\|_F}{\|J_{\theta_B}(p)\|_F}
\]
where $\|\cdot\|_F$ is the Frobenius norm.
\end{definition}

\section{Main Theorem}

\begin{theorem}[IIAR Monotonicity]
For any non-trivial instruction tuning update $\Delta\theta \neq 0$ with bounded perturbation $\|\Delta\theta\|_2 \leq \epsilon$, there exists a prompt $p$ such that
\[
\kappa(\theta_B, \theta_I, p) > 1.
\]
\end{theorem}

\begin{proof}
Let $J_B = J_{\theta_B}(p)$ and $J_I = J_{\theta_I}(p) = J_B + \Delta J$ where $\Delta J = \nabla_p(\Delta\theta \cdot \nabla_\theta h_{\theta_B}(p)) + O(\|\Delta\theta\|^2)$.

By the subadditivity of the Frobenius norm:
\[
\|J_I\|_F = \|J_B + \Delta J\|_F \geq \|J_B\|_F - \|\Delta J\|_F
\]

For non-zero $\Delta\theta$ with bounded perturbation, $\|\Delta J\|_F > 0$ for some prompt $p$ (otherwise $\Delta\theta$ would have zero effect on output, contradicting training objectives).

If $\|J_B\|_F \approx 0$ (base model is insensitive to prompt), then:
\[
\|J_I\|_F \geq \|\Delta J\|_F - \|J_B\|_F \approx \|\Delta J\|_F > 0
\]
giving $\kappa \gg 1$. For non-zero $J_B$, by continuity of the gradient with respect to $\theta$, there exists some $p$ where $J_B$ and $\Delta J$ are not exactly aligned, yielding $\|J_I\|_F > \|J_B\|_F$.
\end{proof}

\begin{corollary}[Differential Direction]
Let $b_f$ and $b_c$ be format bias and content bias. Under IIAR with $\kappa > 1$, there exist positive constants $\alpha, \beta$ such that:
\[
\Delta b_f = -\alpha(\kappa - 1) \quad \text{and} \quad \Delta b_c = +\beta(\kappa - 1)
\]
\end{corollary}

\begin{proof}
Improved attention to format tokens reduces format bias (model better follows rubric). Improved attention to content tokens increases content bias (model more susceptible to exemplar priming). The opposite signs follow because format bias is \emph{failure to follow instructions} (improved attention reduces it) while content bias is \emph{over-attention to distracting exemplars} (improved attention increases it).
\end{proof}

\section{Predictions}

\begin{theorem}[Scaling Law Prediction]
For model size $s$, the bias change $|\Delta b(s)|$ scales as:
\[
|\Delta b(s)| = \gamma \cdot s^\delta + O(1)
\]
with $0 < \delta < 1$, where $\gamma$ is architecture-dependent.
\end{theorem}

\textit{Intuition:} Larger models have more attention capacity, enabling greater redistribution. Returns are diminishing (sub-linear) due to saturation.

\begin{corollary}[Bias Bound]
For any model family with instruction tuning:
\[
|\Delta b_f(rubric)| + |\Delta b_f(score)| \geq |\Delta b_c(ref)|
\]
\end{corollary}

\begin{proof}
By Corollary 1, format bias changes are $-\alpha(\kappa-1)$ and $-\alpha'(\kappa-1)$, content bias change is $+\beta(\kappa-1)$. Since format bias subsumes two independent bias types ($\alpha + \alpha' > \beta$ for typical models), the sum of format reductions exceeds the content increase. Our empirical data confirms this bound: $0.44 + 0.77 = 1.21 \geq 0.35$.
\end{proof}

\section{Future Directions}

The IIAR framework suggests several open questions:

\begin{enumerate}
    \item \textbf{Per-head analysis:} Which attention heads drive format parsing vs content sensitivity? Preliminary work suggests fine-tuning primarily affects a small subset of heads.
    \item \textbf{SFT vs RLHF decomposition:} Separately measuring SFT-only and RLHF-only checkpoints would isolate the training stage responsible for each component of the differential effect.
    \item \textbf{Countermeasure design:} Adversarial training on biased exemplars during instruction tuning could reduce content sensitivity while preserving format gains.
\end{enumerate}

\end{document}
